\begin{studiumworked}{M15. Leer y comprobar una salida CEA · Intermedio}
\textbf{Enunciado.} Use el registro reproducible TP a 2500 K y 1 bar del proyecto. Calcule la constante específica a partir de la masa molar de salida, convierta la capacidad calorífica congelada a SI y compruebe conservación elemental. Explique qué prestación de motor no se puede deducir.

\textbf{Datos.} Salida CEA: $17.42656\,\mathrm{kg/kmol}$; Capacidad congelada: $3.04142\,\mathrm{kJ\,kg^{-1}\,K^{-1}}$.

\textbf{Incógnitas.} R, cp SI y alcance del resultado.

\textbf{Hipótesis.} Equilibrio TP restringido de seis especies. Versiones, base de propiedades y unidades conservadas.

\textbf{Modelo físico.} Conservación elemental y conversión de unidades.
\[R=R_u/\mathcal M,\quad c_{p,SI}=1000c_{p,kJ}\]
\[\sum_i x_i=1,\quad (\sum_i a_{Hi}x_i)/(\sum_i a_{Oi}x_i)=2\]
\textbf{Desarrollo.}
\begin{enumerate}
\item La salida registra $\mathcal M=17.42656$ kg/kmol. Usar $R_u=8314.462618$ J/(kmol K) da $477.1145\,\mathrm{J/(kg K)}$.
\item La capacidad congelada se convierte multiplicando por mil: $3041.42\,\mathrm{J/(kg K)}$. No se usa la de equilibrio como si fuese la misma derivada.
\item Las fracciones del registro suman uno dentro de $10^{-8}$ y el cociente H/O es 2. La verificación se hizo con las fracciones sin redondear y el script conserva la salida original.
\end{enumerate}
\textbf{Operaciones para contrastar el cálculo.}
\begin{enumerate}

\item Magnitud 1.

\[\fitmath{y_{1}=\frac{8314.462618}{17.42655628}\approx 477.1145\,\mathrm{J/(kg\,K)}}\]

\item Magnitud 2.

\[\fitmath{y_{2}=3.041419847\cdot 1000\approx 3041.42\,\mathrm{J/(kg\,K)}}\]

\end{enumerate}

\textbf{Resultados.} $R=477.1145$ J/(kg K); $c_{p,fr}=3041.42$ J/(kg K).

\textbf{Interpretación.} El equilibrio no determina prestaciones de motor sin resolver cámara y expansión.

\textbf{Comprobación.} Fracciones normalizadas y H/O conservado con tolerancias declaradas.

\textbf{Límites.} No es ensayo ni cálculo de Isp real.
\end{studiumworked}
